\section{System Implementation and Validation}

This chapter describes the practical realization of the Plant Up! sensor station. With the sensor selection (Section 2.3.3), firmware pipeline (Section 2.3.4), and decision logic (Chapter 2.4) already established, the focus here is on three things that have not yet been covered: how the components were physically wired and housed, how raw ADC readings were calibrated into meaningful moisture values, and how the assembled system was validated against real plant data.

\subsection{Hardware Assembly}
\begin{itemize}
    \item \textbf{Wiring and pin assignment:}
    \begin{itemize}
        \item GPIO pin mapping table: which ESP32 pin connects to which sensor
        \item Wiring diagram or schematic as figure
    \end{itemize}

    \item \textbf{Enclosure:}
    \begin{itemize}
        \item Custom 3D-printed housing for indoor use
        \item Protection against water splashes, ventilation openings for climate sensor
        \item Photo of the finished sensor station as figure
    \end{itemize}
\end{itemize}

\subsection{Sensor Calibration}
\begin{itemize}
    \item \textbf{Problem:}
    \begin{itemize}
        \item ESP32 ADC produces raw values (0--4095) with no inherent unit
        \item A user cannot interpret ``ADC reads 2340'' without context
    \end{itemize}

    \item \textbf{Two-point calibration:}
    \begin{itemize}
        \item Step 1: Record raw value in air-dry substrate (dry reference)
        \item Step 2: Record raw value in fully saturated substrate (wet reference)
        \item Linear mapping between these two endpoints to a 0--100\% moisture scale
        \item Concrete calibration values from the Plant Up! test setup
    \end{itemize}

    \item \textbf{Limitations:}
    \begin{itemize}
        \item Calibration is substrate-specific (different soil types produce different baselines)
        \item Plant Up! provides relative moisture indication, not laboratory-grade VWC
        \item Recalibration recommended when substrate or pot changes
    \end{itemize}
\end{itemize}

\subsection{Practical Validation}
\begin{itemize}
    \item \textbf{Test setup:}
    \begin{itemize}
        \item Specific plant type, location, and monitoring duration
        \item Controlled watering event during test phase
    \end{itemize}

    \item \textbf{Expected vs. measured behavior:}
    \begin{itemize}
        \item Soil moisture: spike after watering, gradual decline from evaporation and uptake (consistent with theory from Section 2.2)
        \item Light: daily cycle visible (bright daytime, low at night)
        \item Temperature and humidity: stable baseline with minor shifts from heating or ventilation
    \end{itemize}

    \item \textbf{Evidence:}
    \begin{itemize}
        \item At least one figure showing real sensor data from the Plant Up! app
        \item Visual comparison between measured curve and expected physical behavior
    \end{itemize}

    \item \textbf{Limitations:}
    \begin{itemize}
        \item Limited number of test plants and short test duration
        \item No laboratory reference sensor for absolute accuracy verification
        \item Results validate functional correctness, not scientific precision
    \end{itemize}
\end{itemize}
